\documentclass[11pt,letterpaper]{article}
\usepackage[margin=1in,headheight=14pt]{geometry}
\usepackage[T1]{fontenc}
\usepackage{lmodern,amsmath,amssymb,amsthm,mathtools,booktabs,tabularx,longtable,array}
\usepackage{microtype,enumitem,fancyhdr,needspace,xurl}
\usepackage[hidelinks]{hyperref}
\hypersetup{pdfauthor={Cecil Jentges},pdftitle={Ordered-weight certificates for odd cycles in dense coprime graphs}}
\setlength{\parskip}{4pt plus 1pt minus 1pt}
\setlength{\parindent}{1.2em}
\setlength{\emergencystretch}{2em}
\allowdisplaybreaks[1]
\setlist{nosep,leftmargin=1.4em}
\pagestyle{fancy}\fancyhf{}
\fancyhead[L]{\small C. Jentges\quad Ordered-weight certificates}
\fancyhead[R]{\small Mathematical manuscript}
\fancyfoot[C]{\thepage}
\renewcommand{\headrulewidth}{0.3pt}
\numberwithin{equation}{section}
\newtheorem{theorem}{Theorem}[section]
\newtheorem{lemma}[theorem]{Lemma}
\newtheorem{proposition}[theorem]{Proposition}
\newtheorem{corollary}[theorem]{Corollary}
\theoremstyle{definition}\newtheorem{definition}[theorem]{Definition}
\theoremstyle{remark}\newtheorem{remark}[theorem]{Remark}
\newcommand{\NN}{\mathbb N}
\newcommand{\ind}{\mathbf 1}
\newcommand{\pos}[1]{\left(#1\right)_{+}}
\newcommand{\supp}{\operatorname{supp}}
\newcommand{\rad}{\operatorname{rad}}
\newcommand{\wt}{\widetilde\rho}
\newcommand{\abs}[1]{\lvert#1\rvert}
\newcommand{\floor}[1]{\left\lfloor#1\right\rfloor}
\newcommand{\ceil}[1]{\left\lceil#1\right\rceil}
\newcommand{\Oset}{\mathcal O}
\newcommand{\Cset}{\mathcal C}
\newcommand{\Ntilde}{\widetilde N}
\newcommand{\sourcepath}[1]{\path{#1}}
\newcommand{\version}{12 October 2026 (UTC), version 1.1}

\begin{document}
\thispagestyle{plain}
\begin{center}
{\LARGE\bfseries Ordered-weight certificates for odd cycles\\[4pt] in dense coprime graphs\par}
\vspace{9pt}
{\large A computer-assisted argument for Erd\H{o}s Problem 883(i)\par}
\vspace{12pt}
{\large Cecil Jentges\par}
{\small Independent researcher\par}
\vspace{6pt}
{\small\version\par}
\end{center}
\begin{abstract}
For $A\subseteq\{1,\ldots,n\}$, let $G(A)$ join distinct integers precisely when they are coprime. We present a computer-assisted proof that
$|A|>\lfloor n/2\rfloor+\lfloor n/3\rfloor-\lfloor n/6\rfloor$
forces every odd cycle of length at most $\lfloor n/3\rfloor+1$, for every $n$.
The central construction chooses an odd skeleton and fills its gaps with distinct common neighbors. An adjusted totient weight absorbs a finite inclusion--exclusion error into a product bound. Ordering the interior skeleton by this weight bounds the number of deficient connector lists using a single cumulative count, rather than a separate transition count for each divisibility signature. Even connectors and arbitrary unused connectors provide complementary certificates. Explicit finite-prefix distribution bounds cover all $n\ge10^8$; smaller parameters are covered by arithmetic certificates and proof-producing satisfiability computations. The supporting certificate archive is an integral part of the argument. We separate the reusable graph-theoretic mechanism, the arithmetic specialization, the finite evidence, and the independently attributed resolution of the second question.
\end{abstract}
\noindent\textbf{Keywords.} Coprime graphs; odd cycles; ordered Hall condition; totient distribution; computer-assisted proof; exact certificates.

\medskip
\noindent\textit{Review and version note.}
This manuscript consolidates the project proof of 10 October 2026 and its 11 October analytic supplement. The present revision is prepared for mathematical review. It does not assert organizer acceptance, global priority, external human endorsement, or a complete Lean formalization. Substantial AI assistance and the exact scope of checking are described in Section~\ref{sec:evidence}.

\section{The problem and the contribution}\label{sec:statement}
Write $[n]=\{1,\ldots,n\}$ and
\[
 T(n)=\floor{n/2}+\floor{n/3}-\floor{n/6}.
\]
The graph $G(A)$ is simple and undirected. Thus a cycle always means a simple cycle with the stated number of distinct vertices, not merely a closed walk.

\begin{theorem}[All-$n$ odd-cycle assertion]\label{thm:main}
For every integer $n\ge1$ and every $A\subseteq[n]$, if $|A|>T(n)$, then $G(A)$ contains $C_L$ for every odd integer $L$ with
\[
3\le L\le\floor{n/3}+1.
\]
\end{theorem}
The proof combines the mathematical reductions below with the finite certificate assertions in Section~\ref{sec:finite}. Their input files, checks, and coverage are part of the proof, rather than optional experimental support.

The density threshold is sharp. Indeed,
\[
 U_n=\{v\in[n]:2\mid v\text{ or }3\mid v\}
\]
has $T(n)$ elements. Its even vertices form an independent set, as do its odd multiples of $3$. Hence $G(U_n)$ is bipartite and has no odd cycle.

Put $n=6k+r$ with $0\le r\le5$. The largest odd integer at most $\lfloor n/3\rfloor+1$ is $2k+1$. Therefore Theorem~\ref{thm:main} is equivalent to requiring $C_{2\ell+1}$ for every $1\le\ell\le\lfloor n/6\rfloor$. The requested range is empty for $n<6$. There is no unspecified lower bound on $n$ in the theorem.

\subsection{Antecedents and scope}
Erd\H{o}s and S\'ark\"ozy established an odd-cycle result with a positive linear half-length constant and posed the sharp-range problem~\cite{ES97}. S\'ark\"ozy subsequently settled the separate complete-tripartite question~\cite{S99}. A 2026 manuscript of Della Pietra describes an asymptotic route using surplus selection, totient profiles, divisibility signatures, and Hall-type matching~\cite{DP26}. Pan's contemporaneous manuscript presents another all-$n$ completion~\cite{Pan26}.

The contribution developed here is the adjusted-weight pair estimate and its use with an ordered interior skeleton, together with complementary connector families and a finite-to-infinite certificate assembly. The common ingredients---inclusion--exclusion, the totient product, Hall's principle, and extremal density examples---are not claimed as inventions of this work. The methodological distinction concerns their particular combination. This paper does not claim a stronger final odd-cycle theorem than other all-$n$ solutions or assert a settled historical priority.

\subsection{Proof outline}
The proof has four interfaces. First, an ordering argument converts common-neighbor information into distinct connector choices. Second, the adjusted totient profile provides that common-neighbor information. Third, explicit arithmetic counts verify the resulting criteria for every relevant missing-even parameter. Finally, finite certificates cover the range below the analytic threshold. This separation permits a reviewer to check the structural proof without reproducing the entire search that discovered the certificates.

\section{Ordered weights and connector selection}\label{sec:method}
For a real number $z$, write $z_+=\max(0,z)$ and put
\[
 H(c)=c+\ind_{\{c>0\}}\qquad(c\in\NN).
\]
The function $H$ is nondecreasing.

\begin{lemma}[Discarded-tail bound]\label{lem:tail}
Let $Y_0\subseteq Y\subseteq\Omega$ be finite, let $w:\Omega\to\mathbb R$, and suppose
$|Y\setminus Y_0|=\beta$ and $w(y)\le w(y_0)$ for every $y\in Y\setminus Y_0$ and $y_0\in Y_0$. Then for every $t\in\mathbb R$,
\[
\#\{y_0\in Y_0:w(y_0)<t\}
\le\pos{\#\{v\in\Omega:w(v)\le t\}-\beta}.
\]
\end{lemma}
\begin{proof}
If the left side is zero, the claim is immediate. Otherwise choose a retained vertex of weight below $t$. Every discarded vertex has weight below $t$ as well. The discarded vertices and the retained vertices counted on the left are disjoint subsets of the ambient set counted on the right before subtracting $\beta$.
\end{proof}

\begin{lemma}[Ordered-gap bound]\label{lem:gaps}
Let $o_0,\ldots,o_\ell$ be distinct vertices, with $\ell\ge1$, and suppose
\[
 w(o_1)\ge w(o_2)\ge\cdots\ge w(o_{\ell-1}).
\]
The two endpoints are not required to respect this ordering. For a threshold $t$, mark a gap $(o_i,o_{i+1})$ if one of its endpoints has weight below $t$. If $c$ skeleton vertices have weight below $t$, at most $H(c)$ gaps are marked.
\end{lemma}
\begin{proof}
For gaps $i\le\ell-2$, assign a marked gap to $o_{i+1}$ whenever that vertex has weight below $t$. Otherwise assign it to $o_i$. The latter case can occur only for $i=0$: within the nonincreasing middle, a low-weight vertex cannot precede a high-weight vertex. These assignments are injective. The final gap adds at most one. When $c=0$, no gap is marked. The same argument includes $\ell=1$, when the final gap is the only gap.
\end{proof}

\begin{lemma}[Ordered Hall condition]\label{lem:hall}
Finite sets $L_1,\ldots,L_\ell$ have distinct representatives if
\begin{equation}\label{eq:orderedhall}
\#\{i:|L_i|\le j-1\}\le j-1\qquad(1\le j\le\ell).
\end{equation}
\end{lemma}
\begin{proof}
Sort the lists by nondecreasing size. Condition~\eqref{eq:orderedhall} says the $j$th list has at least $j$ members. At the $j$th step at most $j-1$ members have already been chosen, so greedy selection succeeds. This is a sufficient size-based form of Hall's matching principle, not a claim that it is equivalent to Hall's full condition.
\end{proof}

\begin{theorem}[Weighted connector certificate]\label{thm:abstract}
Let $F$ be a finite simple graph. Choose $Y_0\subseteq Y\subseteq\Omega\subseteq V(F)$ as in Lemma~\ref{lem:tail}, with $|Y_0|\ge\ell+1$, and suppose $Y_0$ contains adjacent distinct endpoints $a,a'$. Choose $\ell-1$ other vertices in $Y_0$, order them by nonincreasing weight, and place $a,a'$ at the ends. Write $S=\{o_0,\ldots,o_\ell\}$ for the resulting skeleton. For each gap choose a list
\[
L_i\subseteq\bigl(N_F(o_{i-1})\cap N_F(o_i)\bigr)\setminus S.
\]
Suppose for every $1\le j\le\ell$ there is a threshold $t_j$ such that
\[
\min(w(o_{i-1}),w(o_i))\ge t_j\ \Longrightarrow\ |L_i|\ge j.
\]
Define $B_j=\bigl(\#\{v\in\Omega:w(v)\le t_j\}-\beta\bigr)_+$.
If $H(B_j)\le j-1$ for all $j$, then $F$ contains $C_{2\ell+1}$.
\end{theorem}
\begin{proof}
A list of size at most $j-1$ must correspond to a marked gap at threshold $t_j$. Lemmas~\ref{lem:tail} and~\ref{lem:gaps} bound the number of these gaps by $H(B_j)$. Lemma~\ref{lem:hall} supplies distinct representatives $u_i\in L_i$. The sequence
\[
 o_0,u_1,o_1,u_2,\ldots,u_\ell,o_\ell,o_0
\]
has the required edges; all its vertices except the repeated starting vertex are distinct because the lists avoid $S$. The edge $aa'$ closes the cycle.
\end{proof}

\begin{remark}[What the ordering accomplishes]
A vertex-by-vertex union bound would count up to two affected gaps per low-weight vertex. The ordered middle improves this to $c+1$ when $c>0$, while allowing the coprime endpoints to be prescribed. The first rank still demands zero deficient lists. The certificate deliberately retains that strict requirement rather than replacing it by an average-degree assertion.
\end{remark}

\section{Arithmetic input: the adjusted totient profile}\label{sec:arithmetic}
Fix $n\ge6$ and abbreviate
\[
 q=\floor{n/2},\quad m=\floor{n/6},\quad
 h=\floor{n/3}-\floor{n/6},\quad T=q+h.
\]
It suffices to treat $|A|=T+1$, since larger sets contain a subset of this cardinality. Let
\[
 b=\#\{k\in[q]:2k\notin A\},\qquad Y=\{v\in A:v\text{ odd}\}.
\]
Then $|Y|=h+1+b$, $h\ge m$, and
$0\le b\le\lceil n/2\rceil-h-1$.

\begin{lemma}[Elementary endpoint selection]\label{lem:endpoints}
Every $h+1$ distinct odd integers in $[n]$ contain a coprime pair.
\end{lemma}
\begin{proof}
Place $1$ in group $0$ and each odd $v\ge3$ in group $\lfloor(v-3)/6\rfloor$. The nonunit members of a group are contained in $\{6j+3,6j+5,6j+7\}$ and their pairwise differences are $2$ or $4$. As they are odd, they are pairwise coprime; $1$ is coprime to everything. For odd $v\le n$ the group index is below $h$. Pigeonhole selection proves the claim.
\end{proof}

For a positive integer $v$, let $\omega(v)$ count its distinct prime divisors and put
\[
 \rho(v)=\frac{\phi(v)}v=\prod_{p\mid v}\left(1-\frac1p\right),
 \qquad \wt_x(v)=\rho(v)-\frac{2^{\omega(v)}}{\sqrt{2x}}\quad(x>0).
\]
Products here and below run over distinct primes. The empty product is $1$ and $\omega(1)=0$.

\begin{lemma}[Finite inclusion--exclusion error]\label{lem:legendre}
For a nonempty finite set $P$ of primes and a real $x\ge0$, let $c_P(x)$ count the integers in $[1,x]$ divisible by no prime in $P$. Then
\[
\left|c_P(x)-x\prod_{p\in P}(1-1/p)\right|<2^{|P|-1}.
\]
\end{lemma}
\begin{proof}
For $P=\varnothing$, the error is $\lfloor x\rfloor-x\in(-1,0]$. Adjoining $p\notin P$ gives
$c_{P\cup\{p\}}(x)=c_P(x)-c_P(x/p)$ and the same subtraction for the main terms. The first subtraction gives an error in $(-1,1)$; each additional subtraction doubles the strict symmetric bound.
\end{proof}

\begin{lemma}[Adjusted-product bound]\label{lem:pair}
Let $u,v$ be distinct positive odd integers and let $x$ be a positive integer. If $\wt_x(u),\wt_x(v)\ge0$, then
\begin{equation}\label{eq:pair}
\#\{k\in[x]:\gcd(k,uv)=1\}>x\wt_x(u)\wt_x(v).
\end{equation}
\end{lemma}
\begin{proof}
The prime support of $uv$ is nonempty, since distinct positive integers cannot both be $1$. Put $e=2^{\omega(u)}/\sqrt{2x}$ and $e'=2^{\omega(v)}/\sqrt{2x}$. A prime dividing both endpoints contributes once to $\rho(uv)$ but twice to $\rho(u)\rho(v)$, so $\rho(uv)\ge\rho(u)\rho(v)$. Also
$2^{\omega(uv)-1}\le2^{\omega(u)+\omega(v)-1}=xee'$.
Lemma~\ref{lem:legendre} now gives a strict lower bound by
$x(\rho(u)\rho(v)-ee')$. Finally,
\[
 \rho(u)\rho(v)-ee'-(\rho(u)-e)(\rho(v)-e')
 =e'(\rho(u)-e)+e(\rho(v)-e')\ge0.
\]
\end{proof}
The point of the adjustment is that the pair-dependent error has become a product of two one-vertex weights. No assertion that $u,v$ are coprime to one another is needed in Lemma~\ref{lem:pair}.

\section{Two complementary connector families}\label{sec:families}
Fix $1\le\ell\le m$. For either $w=\wt_q$ or $w=\wt_n$, retain $h+1$ elements $Y_0$ of largest weight in $Y$. Exactly $b$ selected odds are discarded. Lemma~\ref{lem:endpoints} gives coprime endpoints in $Y_0$. Complete them to $\ell+1$ skeleton vertices, ordering the middle by nonincreasing $w$ as in Theorem~\ref{thm:abstract}.

\subsection{Even connectors: family G}
For each gap put
\[
L_i=\{2k\in A:\gcd(k,o_{i-1}o_i)=1\},
\qquad
\Ntilde(y)=\#\{v\le n:v\text{ odd},\ \wt_q(v)\le\sqrt{y/q}\}.
\]
At most $b$ evens are missing. Thus, if both endpoint weights are at least
$t_j=\sqrt{(b+j-1)/q}$, Lemma~\ref{lem:pair} makes the common-even count strictly greater than $b+j-1$, hence at least $b+j$. At least $j$ such evens lie in $A$, and they avoid the odd skeleton. Therefore
\begin{equation}\label{eq:G}
H\bigl((\Ntilde(b+j-1)-b)_+\bigr)\le j-1
\qquad(1\le j\le\ell)
\end{equation}
is a sufficient certificate for $C_{2\ell+1}$. Negative endpoint weights cause no difficulty: the implication is used only when both weights exceed the nonnegative threshold.

\subsection{Arbitrary unused connectors: family GA}
Use $w=\wt_n$ and
\[
 L_i=\{z\in A\setminus S:\gcd(z,o_{i-1}o_i)=1\},\qquad
 N_A(t)=\#\{v\le n:v\text{ odd},\ \wt_n(v)\le t\}.
\]
At most $(n-|A|)+|S|=n-T+\ell$ vertices of $[n]$ are unavailable. Consequently, with
\[
 t_j^A(\ell)=\sqrt{\frac{n-T+\ell+j-1}{n}},
\]
Theorem~\ref{thm:abstract} applies whenever
\begin{equation}\label{eq:GA}
 H\bigl((N_A(t_j^A(\ell))-b)_+\bigr)\le j-1
 \qquad(1\le j\le\ell).
\end{equation}
Different parameter ranges may use different skeletons. Within one application, all rank conditions concern the same skeleton and the same connector lists.

\subsection{A scalar overlap criterion}
Let $y_1=0$ if $\Ntilde(0)>0$; otherwise let
\[
 y_1=\min\{y\in\NN:y\ge1,\ \Ntilde(y)\ge y\}.
\]
This minimum exists: $\wt_q(v)<1$ for every positive $v$, so $\Ntilde(q)=\lceil n/2\rceil\ge q$.
Define
\begin{equation}\label{eq:bA}
b_A=\max\left(0,\ N_A(t_1^A(m)),\
       \max_{2\le j\le m}\{N_A(t_j^A(m))-j+2\}\right),
\end{equation}
where the last maximum is omitted when $m=1$.

\begin{theorem}[G/GA overlap]\label{thm:G}
If $b_A\le y_1-m$, then Theorem~\ref{thm:main} holds at $n$.
\end{theorem}
\begin{proof}
If $b+\ell\le y_1$, then $\Ntilde(b)\le b$ and
$\Ntilde(y)\le y-1$ for $b<y\le b+\ell-1$, so~\eqref{eq:G} holds.
For $j=1$, condition~\eqref{eq:GA} is exactly $N_A(t_1^A(\ell))\le b$. For $j\ge2$, it is equivalent to $N_A(t_j^A(\ell))\le b+j-2$. Thus $b\ge b_A$ gives family GA for $\ell=m$, and also for every smaller $\ell$ because the thresholds decrease. If $b<b_A$, then
$b+\ell\le b_A-1+m\le y_1-1$, and family G applies.
\end{proof}
A convenient stronger condition for finite computations is
\begin{equation}\label{eq:simple}
N_{GA}:=N_A\left(\sqrt{\frac{n-T+2m-1}{n}}\right)\le y_1-m,
\end{equation}
since $b_A\le N_{GA}$. A failure of~\eqref{eq:simple} is a failure of this sufficient certificate, not a counterexample to the cycle theorem.

\subsection{A transparent example}
Take $n=18$ and $A=[13]$, so $T=12$, $q=9$, $h=m=3$, and $b=3$. The four selected odds with largest $\wt_9$ are $1,13,11,7$. Choose skeleton $(1,13,11,7)$: its endpoints are coprime and its interior weights are decreasing. Each gap has the six even connectors $2,4,6,8,10,12$ available. One choice produces
\[
 1,2,13,4,11,6,7,1.
\]
Every edge is coprime and the seven vertices are distinct. This example illustrates the mechanism; the all-subsets assertion follows from the certificates, not from the example.

\section{The infinite range: explicit finite-prefix bounds}\label{sec:analytic}
This section supplies the mathematical reduction of the tail to fixed numerical inequalities. Throughout,
\[
 N_0=10^8,\quad \delta=(N_0-1)^{-1},\quad
 Y_1=\floor{81q/100}.
\]
It is enough to establish
\begin{align}
 \Ntilde(0)&=0, &&\text{(A0)}\label{eq:A0}\\
 \Ntilde(y)&\le y-1 &&(1\le y\le Y_1),\quad\text{(A1)}\label{eq:A1}\\
 b_A&\le Y_1+1-m. &&\text{(A2)}\label{eq:A2}
\end{align}
Indeed, (A0)--(A1) give $y_1\ge Y_1+1$, after which (A2) implies Theorem~\ref{thm:G}. These are inequalities for actual finite prefixes $[n]$, uniformly over $n\ge N_0$. No limiting distribution is used as a replacement for them.

\subsection{Uniform control of the weight adjustment}
Let $\Pi_k$ be the product of the first $k$ odd primes. If $\Pi_k\le n<\Pi_{k+1}$, every odd $v\le n$ has at most $k$ distinct prime divisors and
\[
 \rho(v)\ge g_k:=\prod_{i=1}^k(1-1/p_i).
\]
Thus, writing $e_q(n)=2^k/\sqrt{2q}$ and $e_n(n)=2^k/\sqrt{2n}$,
\begin{align}
 \Ntilde(y)&\le N^0\bigl(\sqrt{y/q}+e_q(n)\bigr),\\
 N_A(t)&\le N^0(t+e_n(n)),
 \qquad N^0(\tau)=\#\{v\le n:v\text{ odd},\rho(v)\le\tau\}.
\end{align}
The quantities $e_q(n)$ and $e_n(n)$ decrease within each primorial interval. At successive odd primorial endpoints $P,pP$,
\[
 \frac{e_q(pP)^2}{e_q(P)^2}=\frac{4(P-1)}{pP-1}<\frac4p<1.
\]
For the normalized shift and reciprocal density scale, the corresponding ratios are
\begin{align}
 \frac{(e_q/g)_{pP}^{\,2}}{(e_q/g)_{P}^{\,2}}
 &=\frac{4p^2(P-1)}{(p-1)^2(pP-1)}<\frac{4p}{(p-1)^2}<1\quad(p\ge7),\\
 \frac{(1/(qg^2))_{pP}}{(1/(qg^2))_{P}}
 &=\frac{p^2(P-1)}{(p-1)^2(pP-1)}<\frac{p}{(p-1)^2}<1.
\end{align}
The first relevant interval is truncated at $N_0$, not at $\Pi_7=4\,849\,845$. Comparing $N_0$ with $\Pi_8=111\,546\,435$ therefore controls every later interval. This proves the infinite continuation of the primorial calculation; a finite list of observed decreases is not being used as an induction argument.

Let $e_q^*=\sup_{n\ge N_0}e_q(n)$ and $e_n^*=\sup_{n\ge N_0}e_n(n)$. They occur at $\Pi_8$ and are respectively
\[
 e_q^*=\frac{256}{\sqrt{111546434}},\qquad
 e_n^*=\frac{256}{\sqrt{223092870}}.
\]
Displayed decimal approximations to these constants are for orientation only. The numerical certificates use outward interval bounds.

\subsection{Small-prime classes}
Let $Q=\{3,5,7,11,13,17,19\}$, put $P=2\prod_{p\in Q}p=9\,699\,690$, and let $\kappa=\phi(P)/P$. The count
$C_P(y)=\#\{1\le w\le y:\gcd(w,P)=1\}$ satisfies
\begin{equation}\label{eq:discrepancy}
 C_P(y)\le\kappa y+D^+,\qquad D^+=\frac{2799708}{323323}.
\end{equation}
The discrepancy is periodic with period $P$. Between consecutive integers it decreases linearly, so its maximum occurs immediately at an integer endpoint. The recorded period scan evaluates the exact integer numerator $P C_P(y)-\phi(P)y$. Its magnitude is below $P^2<2^{63}$, so a signed 64-bit scan suffices. Equation~\eqref{eq:discrepancy} is a finite numerical certificate assertion, with its implementation and output in the source archive~\cite{Archive}.

For $S\subseteq Q$, define
\[
 d_S=\prod_{p\in S}\frac1p\prod_{p\in Q\setminus S}\left(1-\frac1p\right),
 \quad r_S=\prod_{p\in S}\left(1-\frac1p\right),
 \quad U_S=\prod_{p\in S}\frac1{\sqrt p-1}.
\]
Every odd $v$ has a unique decomposition $v=uw$, where the prime support of $u$ is precisely $S=Q\cap\supp(v)$ and $\gcd(w,P)=1$. Let $K_S(n,\tau)$ count the integers in this class with $\rho(v)\le\tau$.

Three bounds are available. First, ignoring all exclusions except divisibility by $\prod_{p\in S}p$ gives
\begin{equation}\label{eq:trivialclass}
 K_S(n,\tau)\le\frac{n}{2\prod_{p\in S}p}+1.
\end{equation}
Second, summing~\eqref{eq:discrepancy} over admissible $u\le n$ gives
\begin{equation}\label{eq:classcount}
 K_S(n,\tau)\le\frac n2d_S+D^+\sqrt n\,U_S.
\end{equation}
For the main term, use $\kappa\sum_{\supp(u)=S}u^{-1}=d_S/2$. For the error, use
$\sum_{u\le n,\supp(u)=S}1\le\sqrt n\sum_{\supp(u)=S}u^{-1/2}=\sqrt n\,U_S$.

For the third bound, let $s>0$ and
\[
 h_s(p)=(1-1/p)^{-s}-1,\quad
 M(s)=\prod_{p>19}\left(1+\frac{h_s(p)}p\right),\quad
 C(s)=\prod_{p>19}\left(1+\frac{h_s(p)}{\sqrt p}\right).
\]
Extend $h_s$ multiplicatively to squarefree integers supported on primes above $19$, with $h_s(1)=1$. Expanding the finite product gives
$\rho(w)^{-s}=\sum_{d\mid w}h_s(d)$.
All terms are nonnegative. For $X\ge1$,
\[
 \sum_{\substack{w\le X\\(w,P)=1}}\rho(w)^{-s}
 \le\kappa X M(s)+D^+\sqrt X\,C(s).
\]
Here the error sum is bounded using $1\le(X/d)^{1/2}$ for $d\le X$. The elementary inequality
$\ind_{\{\rho(w)\le\theta\}}\le\theta^s\rho(w)^{-s}$, followed by summation over $u$, yields
\begin{equation}\label{eq:rankin}
 K_S(n,\tau)\le
 \left(\frac\tau{r_S}\right)^s
 \left(\frac n2d_S M(s)+D^+\sqrt n\,U_S C(s)\right).
\end{equation}
It is useful mainly for $\tau<r_S$, although the displayed bound is valid for every $\tau>0$.

\subsection{One explicit monotone envelope}
Use the fixed exponent set
\[
\mathcal S=\{1,3/2,2,3,4,5,6,8,10,12,15,18,22,26,30,35,40,50,60,80\}.
\]
For $n\ge N_0$, equations~\eqref{eq:trivialclass}--\eqref{eq:rankin} imply
\begin{equation}\label{eq:envelope}
 N^0(\tau)\le\frac n2\Phi(\tau),
\end{equation}
where $\Phi$ is the sum, over $S\subseteq Q$, of the minimum of
\begin{align*}
 &\frac1{\prod_{p\in S}p}+\frac2{N_0},\qquad
 d_S+\frac{2D^+U_S}{\sqrt{N_0}},\\
 &\min_{s\in\mathcal S}\left(\frac\tau{r_S}\right)^s
       \left(d_SM(s)+\frac{2D^+U_SC(s)}{\sqrt{N_0}}\right).
\end{align*}
Each candidate bound is nondecreasing in $\tau$, hence so is $\Phi$.

The infinite Euler products have explicit tails. Put $X=10^7$. For $0<s\le80$ and $p>X$, the inequalities
$-\log(1-z)\le z/(1-z)$ and $\exp(t)-1\le t/(1-t)$ imply
\[
 h_s(p)\le\frac{s}{p-1-s}\le\frac{s}{p-81}<\frac{1.001s}{p}.
\]
Consequently, for $\sigma'=1$ or $1/2$,
\begin{equation}\label{eq:producttail}
 \log\prod_{p>X}\left(1+\frac{h_s(p)}{p^{\sigma'}}\right)
 \le \frac{1.001s}{\sigma'X^{\sigma'}}.
\end{equation}
The last step replaces primes by all integers and bounds their decreasing tail by an integral. The products through $10^5$ are evaluated in interval arithmetic; from $10^5$ to $10^7$, exact prime counts in buckets and monotonicity bound the summands. No unproved assumption about the omitted primes is required.

\subsection{The high part of (A1)}
For $u=\sqrt{y/q}$, we have $n/(2q)\le1+\delta$ and $1/q\le2\delta$. On every rational cell $[u_a,u_b]\subseteq[33/100,9/10]$, the inequality
\begin{equation}\label{eq:highgrid}
 (1+\delta)\Phi(u_b+e_q^*)+2\delta<u_a^2
\end{equation}
implies $\Ntilde(y)\le y-1$ for every represented integer $y$, not just for a sampled endpoint. Indeed, monotonicity bounds $\Ntilde(y)/q$ by the first term, while $1/q\le2\delta$ supplies the unit reserve. The archived interval run checks all $1\,140$ cells of width $1/2000$.

\subsection{The low part of (A1) and (A0)}
Use the same moment expansion over all odd primes, with exponent $4$. If $M_o,C_o$ are the corresponding two Euler products, then
\[
 N^0(\tau)\le\tau^4\left(\frac n2M_o+\sqrt n\,C_o\right).
\]
Write $g=g_k$ on the relevant primorial interval. If $\tau<g$, the count is zero. For $\tau=u+e_q(n)\ge g$, it suffices to check
\begin{multline}\label{eq:lowrankin}
 (33/100+e_q^*)^2(1+\delta)
 \left(M_o+\frac{2C_o}{\sqrt{N_0}}\right)
 +\sup_{n\ge N_0}\frac1{qg^2}\\
 <\left(1-\sup_{n\ge N_0}\frac{e_q(n)}g\right)^2.
\end{multline}
To see this, multiply the corresponding inequality with $\tau^2$ in the first factor by $\tau^2$. Since $\tau\ge g$ and $e_q<g$,
\[
 \frac{\Ntilde(y)}q+\frac1q
 \le \tau^4(1+\delta)(M_o+2C_o/\sqrt{N_0})+\frac1q
 <(\tau-e_q)^2=u^2.
\]
The primorial comparisons above control both suprema in~\eqref{eq:lowrankin}. They also give $e_q(n)<g_k$, so no adjusted weight is nonpositive; this proves (A0).

\subsection{The overlap bound (A2)}
Set $\xi=(j-1)/n\in[0,1/6]$. The exact floor estimates
\[
 T(n)\ge\frac{2n}3-\frac76,\quad
 m\le\frac q3+\frac16,\quad j-2\ge2q\xi-1
\]
give
\[
 (t_j^A(m))^2\le\frac12+\frac7{6N_0}+\xi,
 \qquad Y_1+1-m\ge\left(\frac{81}{100}-\frac13\right)q-\frac16.
\]
It is sufficient on each rational cell $[\xi_a,\xi_b]\subseteq[0,1/6]$ that
\begin{multline}\label{eq:GAgrid}
 (1+\delta)\Phi\left(\sqrt{\frac12+\frac7{6N_0}+\xi_b}+e_n^*\right)
 -2\xi_a+3\delta\\
 <\frac{81}{100}-\frac13.
\end{multline}
Indeed, $3\delta\ge(7/6)/q$ pays for the additive $1$ in the estimate for $N_A(t_j^A(m))-j+2$ and the $1/6$ in the lower bound for $Y_1+1-m$. For $j=1$, this establishes the stronger bound on $N_A(t_1^A(m))+1$, hence also the term required in~\eqref{eq:bA}. The archived run covers the interval by $400$ cells.

\subsection{Numerical certificate interface and conclusion}
The existing 11 October 2026 replay of \sourcepath{analytic_interval.py} reports success for all three uniform bounds. Its displayed margins are approximately
\[
0.0245451344\quad\text{(high A1)},\qquad
0.3646407885\quad\text{(low A1)},\qquad
0.0589644029\quad\text{(A2)}.
\]
In particular the last margin is \emph{greater than $0.058$}, not asserted to be at least $0.059$. These displayed numbers summarize the archived interval output; this typesetting revision did not rerun that entire calculation. The mathematical supplement records the primorial continuation, Euler-product tail, period scan, and grid-to-uniform deductions~\cite{TailReplay}.

Subject to these explicitly identified numerical certificates, (A0)--(A2) hold for every $n\ge10^8$. Together with the overlap theorem, they establish the infinite portion of Theorem~\ref{thm:main}. The smaller values of $n$ are addressed next.

\section{Finite certificates and the all-$n$ assembly}\label{sec:finite}
The finite range is not inferred from samples. The source archive records complete parameter scans, a finite list of residual satisfiability obligations, and interval certificates that dominate every intermediate integer. Appendix~\ref{app:small} consolidates the small-parameter reductions without retaining superseded draft formulas.

\subsection{A coverage map}
\begin{center}
\small
\begin{tabularx}{\linewidth}{@{}p{0.22\linewidth}X@{}}
\toprule
Range & Proof component and evidence interface\\
\midrule
$1\le n\le5$ & The requested odd-cycle range is empty.\\[3pt]
$6\le n<4000$ & The E/R/M arithmetic certificate families, with the $99$ specified residual SAT instances for $n\le79$.\\[3pt]
$4000\le n\le30000$ & The stronger G/GA criterion~\eqref{eq:simple}, except at $4025,4029,4035$, which are covered by the small-parameter certificate families.\\[3pt]
$30001\le n\le10^8$ & $167$ interval/block certificates for G/GA; the last block ends at $10^8$.\\[3pt]
$n\ge10^8$ & The uniform analytic inequalities (A0)--(A2). The overlap at $10^8$ is intentional.\\
\bottomrule
\end{tabularx}
\end{center}
These are the coverage assertions of the frozen project evidence, not a claim that every listed computation was rerun for this revision. The hash-bound ledger supplies the machine-readable inventory and confirms that all required records are present~\cite{Archive}.

\subsection{Dominating an entire block}
For a block $[N_-,N_+]$, use all odd $v\le N_+$, $q_-=\lfloor N_-/2\rfloor$, and $m_+=\lfloor N_+/6\rfloor$. Compute G weights at $q_-$ and GA weights at $N_-$. For the latter use threshold square
\[
 R=\frac13+\frac7{6N_-}+\frac{2m_+-1}{N_-}.
\]
For every $n$ in the block, the ambient odd universe can only shrink; the adjustments $2^{\omega(v)}/\sqrt{2q}$ and $2^{\omega(v)}/\sqrt{2n}$ can only decrease; and the displayed threshold dominates $(n-T(n)+2m(n)-1)/n$. Therefore the block G cumulative counts dominate each per-$n$ count and the block GA count dominates each per-$n$ GA count. Its first-failure value can only be smaller, and $m_+\ge m(n)$. A successful block certificate consequently proves~\eqref{eq:simple} at every integer in the block.

\subsection{The residual satisfiability obligations}
After applying the arithmetic families, the recorded $86$ triples $(n,b,\text{first uncovered }\ell)$ expand to $99$ required target instances. Each asks whether a set $A\subseteq[n]$ of size $T+1$, with exactly $b$ missing evens, avoids a specified odd cycle. The reduction to a finite SAT instance has the following mathematical ingredients.

\paragraph{Support compression.}
If $x\in A$, $y\notin A$, $\supp(x)\subsetneq\supp(y)$, and $x,y$ have the same parity, replace $x$ by $y$. Any new cycle using $y$ can be converted to one in the old set by replacing $y$ by $x$: coprimality with $y$ implies coprimality with $x$. Cardinality and $b$ are unchanged. The total number of prime factors in the selected supports strictly increases, so repeated compression terminates.

\paragraph{Twin classes.}
Integers with equal nonempty prime support are nonadjacent twins in the coprime graph. Permuting them preserves the graph. The empty support occurs only at $1$. Consequently a putative counterexample may be represented by a prefix in each support class. This is a reduction of all counterexamples, not a claim that arbitrary sets themselves have the canonical form.

\paragraph{The additional six instances.}
In six cases the archive also restricts $c$, the number of absent odd multiples of $3$, to $c\le c_{\max}$. Compression cannot increase $c$: when an odd support gains the prime $3$, the replacement can only decrease the number of missing multiples of $3$. All excluded $c>c_{\max}$ cases are separately covered by family M. The further prime-swap symmetries exchange $kp$ with $kp'$ when $3<p<p'$, $p^2>n$, and $\lfloor n/p\rfloor=\lfloor n/p'\rfloor$. Their multipliers are smaller than both primes, so the swaps preserve coprimality, parity, and divisibility by $3$.

\paragraph{Checkable cycle exclusions and UNSAT proofs.}
A discovered cycle justifies excluding any canonical set containing its vertex pattern. The certificate includes the canonical cycle witnessing each such exclusion. The source records $670\,834$ such witnesses across the $99$ instances. A separate encoder reconstructs the CNF, checks each witness and structural constraint, and supplies an LRAT proof of unsatisfiability. The recorded checks use \sourcepath{drat-trim} and \sourcepath{lrat-check}. A heuristic cycle search need not be complete: failure to find a cycle cannot certify an instance; successful completion requires a checked UNSAT certificate for the final sound encoding. The size-$T$ negative control returns the extremal set $U_n$.

\begin{proof}[Assembly of Theorem~\ref{thm:main}]
Reduce to $|A|=T+1$ and write the requested length as $2\ell+1$ with $1\le\ell\le m$. For $n<6$ there is no requested length. Below the analytic threshold, the coverage ledger assigns the relevant $(n,b,\ell)$ either a valid arithmetic certificate or a checked residual unsatisfiability proof. Appendix~\ref{app:small}, Section~\ref{sec:families}, and the counterexample-preserving reductions above show that these certificates imply the required cycle. The block domination argument handles every intermediate integer, rather than only block endpoints. Section~\ref{sec:analytic} covers every remaining $n$. Finally, any cycle in the chosen size-$T+1$ subset also lies in the original set.
\end{proof}

\section{Coverage of the complete original problem}\label{sec:both}
The newly claimed contribution of this manuscript is question (i). Question (ii) asks whether, for every fixed positive integer $\ell$, there is a threshold $N(\ell)$ such that all $n\ge N(\ell)$ and all $A\subseteq[n]$ with $|A|>T(n)$ contain a not necessarily induced $K_{1,\ell,\ell}$.

This second conclusion is the established theorem of G\'abor N. S\'ark\"ozy~\cite{S99}. Its complete-tripartite conclusion has part sizes tending to infinity with $n$, under the same strict density threshold; restricting the two large parts to any fixed $\ell$ gives the requested statement. In the original notation, $f(n,2)$ is the number of integers divisible by $2$ or $3$, which equals $T(n)$ by inclusion--exclusion. The threshold depends on $\ell$, not on the particular set $A$. Complete tripartite subgraph containment permits extra edges within the two large parts.

Thus the proposed original-problem coverage consists of Theorem~\ref{thm:main} together with this cited historical result. No new-solver credit for S\'ark\"ozy's conclusion is claimed. This separates full mathematical coverage from the scope of the present author's contribution; it does not purport to settle an organizer's attribution decision.

\section{Evidence, reproducibility, and responsibility}\label{sec:evidence}
\subsection{Frozen source and verification boundaries}
The mathematical source is the \sourcepath{solver/jsp000734/hall_tail_20261010} directory in the project repository, retained at commit
\begin{center}\small\texttt{a04e7031546461920d8a3ac3f920447f1e689feb}.\end{center}
The earlier mathematical snapshot is
\texttt{66321f74ed35db1df594fd026ae3b2c50bc9528d}.
The 11 October supplement is stored separately in
\sourcepath{solver/jsp000734/submission_math_only_20261011}~\cite{TailReplay}.
The original files and manifest should remain unchanged; this paper supersedes their exposition where the revision record says so, not their historical identity.

A public mathematical submission must expose the indispensable manuscript, data, checkers, and version identifiers. A private repository URL is not a public certificate archive. The accompanying release-assembly tool exports the selected source and its manifest dependencies without publishing unrelated contest work. The PR-preparation tool inserts public links only after the release has been assembled and its public commit specified.

\subsection{Checking recorded in the original development}
The project records several AI-assisted adversarial reviews and code-disjoint reimplementations. Such records are evidence about what was checked and by which program, not independent human mathematical endorsement. The hash-bound ledger \sourcepath{code/coverage_ledger2.py} checks the expected archive. Its deep mode reconstructs SAT inputs and rechecks LRAT proofs; the scanner mode reruns the E/R/M calculations. The historical G/GA implementation uses floating-point filtering followed by an exact predicate at uncertain boundaries. It must not be described as literally integer-only at every operation. Appendix~\ref{app:arithmetic} gives the exact predicate, a bounded-error justification for the historical filter under stated arithmetic assumptions, and a new integer-only alternative.

\subsection{Checks actually performed for this revision}
The accompanying new \sourcepath{verification/integer_replay.py} uses directed rational intervals and exact integer fallback. Its bounded tests checked $150\,000$ least-threshold rows, $500\,000$ count comparisons against the exact predicate, $10\,000$ wide-product decompositions, and $58\,471$ inclusive-equality configurations. The combinatorial tests covered $20\,200$ ordered-gap configurations, $4\,108$ reserve/deletion configurations, and the worked example. A further complete integer-only replay checked every $n$ from $4000$ through $30000$: $25\,998$ cases passed the stronger G/GA certificate and exactly $4025,4029,4035$ failed it, as expected. Those three cases retain their separate E/R/M coverage. Exact details appear in \sourcepath{verification/TEST_RESULTS.json} and \sourcepath{verification/COMPLETE_PER_N_REPLAY.json}; every per-$n$ row is supplied.

The per-$n$ replay is complete on $[4000,30000]$, but this revision does not claim a new replay of all $167$ blocks, the E/R/M scans, all $99$ LRAT certificates, or the analytic interval calculation. The complete historical evidence remains necessary.

\subsection{Authorship, AI assistance, and scope of the claims}
Cecil Jentges selected and directed the research and requested its development, auditing, and submission. AI systems substantially assisted mathematical development, programming, checking, and exposition. This disclosure is not a claim that every proof line was independently checked by the human author. No independent external human verifier is identified in this revision.

The earlier proof development, the later comparison with related work, and any subsequent changes are distinguished in the accompanying provenance record. This revision does not assert that earlier sessions never encountered a related source. Independence and priority require supporting history, not a statement of intent or a newly assigned document date. A later comparison does not by itself erase an earlier contribution, but improvements borrowed after that comparison must be credited.

\appendix
\section{Small-parameter arithmetic certificates}\label{app:small}
This appendix gives the consolidated E/R/M reductions used in the finite archive. It incorporates the corrected single-gap rule, the one-per-later-endpoint count for sorted chunks, and the corrected skeleton-pool size for family M. Historical preliminary bounds are not used as the current statements.

\subsection{Notation and elementary bounds}
Continue with $q,m,h,T,b,Y$ from the main text. Put
\[
 \Delta=\floor{q/3}=m,\qquad \beta_0=q-\Delta,
 \quad\Oset_3=\{v\in[n]:v\text{ odd},\ 3\mid v\},
 \quad\Cset=\{v\in[n]:\gcd(v,6)=1\}.
\]
Let $X=A\cap\Cset$ and $c=|\Oset_3\setminus A|$. Then $|X|=b+c+1$.
Choose any finite set $P$ of signature primes at least $5$. For odd $v$, write
\[
 S(v)=P\cap\supp(v),\qquad
 L(v)=\supp(v)\setminus(\{2,3\}\cup P).
\]
The key of $v$ is $(\ind_{\{3\mid v\}},S(v))$. Let $K$ and $K_C$ count the keys realized among all odds and among $\Cset$, respectively.

To avoid confusing counting losses with the totient density $\rho$, denote the exact losses by
\begin{align*}
 d(v)&=\#\{k\le q:3\nmid k,\ \gcd(k,v)>1\},\\
 d_N(v)&=\#\{u\le n:\gcd(u,v)>1\}.
\end{align*}
For $S\subseteq P$, define the signature losses $\sigma_3(S),\sigma_A(S),\sigma_N(S)$ by counting, respectively, noncoprime integers against $\prod_{p\in S}p$ in the pools $\{k\le q:3\nmid k\}$, $[q]$, and $[n]$. Define
\begin{align*}
 \lambda_3(v)&=\sum_{p\in L(v)}\#\{k\le q:3\nmid k,\ p\mid k\},\\
 \lambda_A(v)&=\sum_{p\in L(v)}\floor{q/p},\qquad
 \lambda_N(v)=\sum_{p\in L(v)}\floor{n/p}.
\end{align*}
All these quantities are integers. For a discount $0\le\eta\le\Delta$, the costs are
\[
 c_3(v)=\sigma_3(S(v))+2\lambda_3(v),\qquad
 c_\eta(v)=\sigma_A(S(v))+2\lambda_A(v)-\eta,
 \qquad c_R(v)=\sigma_N(S(v))+2\lambda_N(v).
\]
The first is used on $\Oset_3$, and the other two on $\Cset$.

The set $\Cset$ partitions into $\{1\}$ and the nonempty intersections with blocks
$\{12j+5,12j+7,12j+11,12j+13\}$, $j\ge0$. Each block is pairwise coprime: a common prime would divide a difference in $\{2,4,6,8\}$, but its members are coprime to $6$. If the number of groups is $g_C(n)$, any $g_C(n)+1$ elements of $\Cset$ contain a coprime pair.

\subsection{Ambient rough-count bounds}
Suppose a retained pool consists of the $s$ smallest-cost elements of a selected set, with at least $\beta$ selected elements discarded. For an ambient subset $W$, let
\[
 \tau(W)=\min_{v\in W}c(v),\qquad
 R(t)=\#\{v\text{ in the ambient universe}:c(v)\ge t\}.
\]
For $W=\varnothing$, define its rough bound to be zero. Otherwise set
\[
 \langle W\rangle_{c,\beta}
 =\min\bigl(|W|,(R(\tau(W))-\beta)_+\bigr).
\]
The number of retained vertices in $W$ is at most this bound. Indeed, if a retained vertex lies in $W$, all discarded selected vertices have cost at least $\tau(W)$, so the same disjoint-count argument as in Lemma~\ref{lem:tail} applies. The absolute bound $|W|$ is independently valid.

\subsection{Family E: even connectors}
Use the $h+1$ smallest-cost selected odds, with cost $c_3$ on $\Oset_3$ and $c_\eta$ on $\Cset$. The subtraction is $\beta=b$. Choose coprime endpoints from the retained pool.
For $\ell=1$, a sufficient separate rule is
\begin{equation}\label{eq:Eone}
 \left\langle\left\{v:d(v)\ge\ceil{(\beta_0-b)/2}\right\}\right\rangle=0.
\end{equation}
It gives $d(a)+d(a')<\beta_0-b$, leaving at least one common available even connector.

For $\ell\ge2$, choose a separator level $g\ge0$. Require
\[
 h+1-\langle\{v:d(v)\ge g+1\}\rangle\ge K+3.
\]
After excluding the two endpoints there are enough separators to choose
$\min(K+1,\ell-1)$ dividers. If $\ell-1\le K+1$, the middle consists entirely of dividers. Otherwise divide the remaining middle vertices into their at most $K$ key classes, order each chunk by increasing cost, and interpose dividers before, between, and after the chunks. Unused dividers stand consecutively.

Let $\theta_j=\beta_0-b-j+1$. The rank-$j$ sufficient condition is
\begin{multline}\label{eq:Ecert}
 \left\langle\{v\in\Oset_3:c_3(v)\ge\theta_j\}
       \cup\{v\in\Cset:c_\eta(v)\ge\theta_j+\Delta-\eta\}\right\rangle\\
 +\min\bigl(2K+2,2\langle\{v:d(v)\ge\theta_j-g\}\rangle\bigr)
 +K\ind_{\{\theta_j\le2g\}}\le j-1
 \qquad(1\le j\le\ell).
\end{multline}
For a same-key $\Oset_3$ gap, the common-even list has size at least
$\beta_0-b-\sigma_3(S)-\lambda_3(v)-\lambda_3(v')$.
For a same-key $\Cset$ gap, its size is at least
$q-b-\sigma_A(S)-\lambda_A(v)-\lambda_A(v')$.
A bad within-chunk gap is therefore charged to its later, larger-cost endpoint. Each vertex receives at most one such charge. A separator-adjacent list has size at least $\beta_0-b-g-d(v)$; there are at most $2K+2$ such gaps and each rough endpoint participates at most twice. A divider--divider gap has size at least $\beta_0-b-2g$, and there are at most $K$ of them. These three counts give~\eqref{eq:Ecert}, then Lemma~\ref{lem:hall} gives the cycle.

\subsection{Family R: a skeleton in $\Cset$}
Put $s_0=\max(\ell+1,g_C+1)$. If $b+1\ge s_0$, retain the $s_0$ smallest-$c_R$ elements of $X$; the valid uniform subtraction is $\beta=b+1-s_0$. The pool contains coprime endpoints. For $\ell=1$, it suffices that
\[
 \left\langle\left\{v\in\Cset:d_N(v)\ge\ceil{(T-1)/2}\right\}\right\rangle=0.
\]
For $\ell\ge2$, choose separators with $d_N\le g_N$ and require
\[
 s_0-\langle\{v\in\Cset:d_N(v)\ge g_N+1\}\rangle\ge K_C+3.
\]
Use the same chunk/divider arrangement, now inside $\Cset$. Connector lists lie in $A\setminus S$, a pool of size $T-\ell$. With $\theta_j^R=T-\ell-j+1$, the condition is
\begin{multline}\label{eq:Rcert}
 \langle\{v\in\Cset:c_R(v)\ge\theta_j^R\}\rangle
 +\min\bigl(2K_C+2,2\langle\{v\in\Cset:d_N(v)\ge\theta_j^R-g_N\}\rangle\bigr)\\
 +K_C\ind_{\{\theta_j^R\le2g_N\}}\le j-1.
\end{multline}
The common-neighbor loss within a key is at most
$\sigma_N(S)+\lambda_N(v)+\lambda_N(v')$; at a separator it is at most $g_N+d_N(v)$. Sorted chunks and the same gap classification prove~\eqref{eq:Rcert} sufficient.

\subsection{Family M: a mixed skeleton}
This family is checked for every possible $c\in[0,h]$. If $|X|=b+c+1\ge s_0$, apply family R with the sharper subtraction $|X|-s_0$. The condition is $|X|\ge\max(\ell+1,g_C+1)$, not merely $|X|\ge\ell+1$.
If $|X|<g_C+1$, the mixed family does not apply. Otherwise assume
\[
 g_C+1\le |X|\le\ell,\qquad z=\ell+1-|X|\ge1.
\]
Put all of $X$ and the $z$ smallest-$c_3$ selected odd multiples of $3$ into the skeleton. The latter pool has subtraction $\beta_3=h-c-z\ge0$. Choose the coprime endpoints inside $X$. Require enough $\Cset$ separators:
\[
 |X|-\#\{v\in\Cset:d_N(v)\ge g_N+1\}\ge K_C+3+z.
\]
Use $K_C+1+z$ separators as dividers. Arrange the remaining $\Cset$ vertices in increasing-cost key chunks, and put each of the $z$ multiples of $3$ between two dividers. Further separators belong to their ordinary key chunks. Let
\[
 g_d=\max\{d(v):v\in\Cset,\ d_N(v)\le g_N\}.
\]
The availability condition ensures this maximum is over a nonempty set. On a $\Cset$--$\Cset$ gap use all unused vertices as possible connectors; on a separator--$\Oset_3$ gap use $3$-free even connectors. The latter list has size at least $\beta_0-b-g_d-d(v)$.

With the same $\theta_j,\theta_j^R$, a sufficient rank bound is
\begin{align}\label{eq:Mcert}
 &\#\{v\in\Cset:c_R(v)\ge\theta_j^R\}\nonumber\\
 &\quad+\min\left(2K_C+2,2\#\{v\in\Cset:d_N(v)\ge\theta_j^R-g_N\}\right)\nonumber\\
 &\quad+(K_C+z)\ind_{\{\theta_j^R\le2g_N\}}
 +\min\left(2z,2\langle\{v\in\Oset_3:d(v)\ge\theta_j-g_d\}\rangle_{c_3,\beta_3}\right)
 \le j-1.
\end{align}
The four terms count within-chunk gaps, separator-adjacent $\Cset$ gaps, consecutive-divider gaps, and the $2z$ gaps incident to the inserted multiples of $3$. Applying ordered Hall proves the assertion.

\subsection{Finite quantifiers}
For each $n$, every admissible $b$ and every $1\le\ell\le m$ must be certified. A chosen family and its parameters may depend on $(n,b,\ell)$; a family-M certificate must account for every $c$. When a case is not certified, a separate residual certificate is required. The negative outcome of one sufficient family never closes a case. The source scanners recompute keys and costs for each allowed signature-prime set~\cite{Archive}.

\section{Exact arithmetic and a fully integer verification route}\label{app:arithmetic}
\subsection{An exact comparison predicate}
Let $f,r,x,Q$ be positive integers, $0<f\le r$, $w\ge0$, and $P\ge0$. Define
\[
 X=2xQf^2-2xPr^2-4^wQr^2.
\]
Then
\begin{equation}\label{eq:exactpredicate}
 \frac fr-\frac{2^w}{\sqrt{2x}}\le\sqrt{P/Q}
 \quad\Longleftrightarrow\quad
 X\le0\ \text{or}\ X^2\le4^{w+1}r^4(2x)PQ.
\end{equation}
Indeed, move the adjustment to the right and square two nonnegative sides. The resulting inequality is
$X\le2^{w+1}r^2\sqrt{2xPQ}$. When $X\le0$ it is automatic. Otherwise squaring preserves equivalence and gives~\eqref{eq:exactpredicate}. The formula includes $P=0$ and equality cases. For the totient profile, use $r=\rad(v)$ and $f=\phi(r)$.

\subsection{What is required to justify the historical floating filter}
The historical checker does not evaluate~\eqref{eq:exactpredicate} on every row. It uses binary64 filtering and exact fallback near a boundary. A precise sufficient execution model is round-to-nearest binary64 arithmetic for division, multiplication, subtraction and square root; exact conversion of the bounded integers; and exact output of $2^w$ for $w=0,\ldots,8$. Finite values and the absence of unaccounted approximate/fast-math operations are also required. These are numerical execution assumptions, not new mathematical hypotheses on the graph.

For the archived domain $N,x\le10^8$, we have $0<f\le r\le10^8$, $w\le8$, and all converted integers are exactly representable. Assume $0\le P/Q\le4$. Let $u=2^{-53}$, $G=f/r$, $e=2^w/\sqrt{2x}$, $z=G-e$, and $t=\sqrt{P/Q}$. Direct propagation of the above rounding bounds gives
\[
 |\widehat G-G|\le u,\qquad
 |\widehat e-e|<3eu,\qquad
 |\widehat z-z|<1100u.
\]
Here $G\le1$, $e<256$, and $t\le2$ are deliberately loose bounds. The two operations defining $\widehat t$ contribute less than $4u$ absolute error. The final subtraction therefore gives
\[
 |\operatorname{fl}(\widehat z-\widehat t)-(z-t)|<1400u<1.6\cdot10^{-13}.
\]
Thus a $10^{-9}$ sign guard is sufficient under the stated execution model.

For the least-threshold computation, if either $z$ or $\widehat z$ is positive, the preceding global estimate implies $e<2$. Repeating the estimate locally gives $|\widehat z-z|<12u$. The map $z\mapsto\max(z,0)$ is $1$-Lipschitz, and its true value is at most $1$. Squaring and multiplying by $x$ then give
\[
 \left|\operatorname{fl}\bigl(x\max(\widehat z,0)^2\bigr)-x\max(z,0)^2\right|<64ux.
\]
If both are nonpositive, both clipped values are zero. A guard of $10^{-12}x+10^{-9}$ around an integer boundary exceeds this error, including the rounding of the guard and fractional-part calculation. Away from that boundary the ceiling is unchanged; uncertain rows must use~\eqref{eq:exactpredicate}. Exact comparisons also correct the sign boundary at zero.

This supplies an explicit error argument for the stated execution model. It does not retroactively certify an undocumented software environment or turn the historical implementation into an integer-only one. The next route eliminates the floating-point premises entirely.

\subsection{Directed rational bounds without floating point}
The new verifier uses $S=2^{32}$. Obtain bounds on the scaled profile using
\[
 g_-=\floor{Sf/r},\qquad g_+=g_-+1.
\]
Put $a=\lfloor S\sqrt{2x}\rfloor=\operatorname{isqrt}(2xS^2)$ and
\[
 e_-=\floor{2^wS^2/(a+1)},\qquad
 e_+=\ceil{2^wS^2/a}.
\]
Then $(g_--e_+)/S\le\wt_x(v)\le(g_+-e_-)/S$. Similarly,
\[
 t_-=\operatorname{isqrt}\bigl(\floor{PS^2/Q}\bigr),\qquad t_+=t_-+1
\]
bound $S\sqrt{P/Q}$. Accept only when the upper weight bound is at most $t_-$, reject only when the lower weight bound exceeds $t_+$, and use~\eqref{eq:exactpredicate} otherwise. Every decision is certified by integers; floating point is not used even to select a side of the comparison.

For
\[
 k_x(v)=\min\{k\ge0:\wt_x(v)\le\sqrt{k/x}\}
       =\ceil{x\max(\wt_x(v),0)^2},
\]
apply the same directed weight bounds to obtain lower and upper ceilings. If they agree, the common value is certified. Otherwise use exact binary search between them.

\subsection{Overflow bounds}
The vectorized implementation uses unsigned 64-bit integers for $Sf$ and for the square of a scaled positive weight $W<S$. Since $f\le10^8$, $Sf<2^{59}$; also $W^2<2^{64}$. To avoid forming the potentially larger product $xW^2$, write $W^2=z_0+Sz_1$ and compute
\[
 p_0=xz_0,\qquad p_1=xz_1+\floor{p_0/S}.
\]
Then $\lfloor xW^2/S^2\rfloor=\lfloor p_1/S\rfloor$. Its remainder is nonzero precisely when either $p_0\bmod S$ or $p_1\bmod S$ is nonzero. Both products stay below $2^{60}$ on the enforced domain. Signed weight differences have magnitude below $2^{40}$. The exact fallback uses Python's arbitrary-precision integers. Histogram counts and cumulative counts are at most the ambient number of odds and fit signed 64-bit integers.

These bounds justify the new implementation's arithmetic structure; the accompanying tests are additional implementation checks, not substitutes for these bounds.

\section{Methodological consequences and boundaries}\label{app:methodconsequences}
\begin{corollary}[Connector reserve]
In Theorem~\ref{thm:abstract}, replace the good-list conclusion $|L_i|\ge j$ by $|L_i|\ge j+r$, where $r\ge0$ is an integer, and retain $H(B_j)\le j-1$. After deletion of any at most $r$ candidate connector vertices, the same skeleton can still be completed to a cycle.
\end{corollary}
\begin{proof}
The ordered list sizes satisfy $d_j\ge j+r$. Deletion of at most $r$ vertices leaves the $j$th list, in that original order, with at least $j$ choices before earlier representatives are removed. Greedy selection succeeds.
\end{proof}
\begin{corollary}[Counting labelled completions]
For fixed ordered skeleton and connector lists whose nondecreasing sizes are $d_1,\ldots,d_\ell$ with $d_j\ge j$, the number of injective choices of representatives is at least
\[
 \prod_{j=1}^{\ell}(d_j-j+1).
\]
In particular $d_j\ge j+r$ gives at least $(r+1)^\ell$ labelled connector assignments.
\end{corollary}
\begin{proof}
At the $j$th greedy step, no more than $j-1$ candidates have been used. Multiply the lower bounds over the decision tree. This counts assignments to the fixed lists; it is not an unqualified enumeration of unlabelled cycles across different skeletons.
\end{proof}
These are elementary consequences of the construction, not a claim of new global resilience or counting bounds for every dense coprime graph. Such stronger applications require the corresponding reserve inequalities, which are not established here. The separate methods note explains the abstraction and its possible research directions; it is a companion exposition rather than a second independently claimed resolution of Problem 883.

\begin{thebibliography}{9}
\bibitem{ES97}
P. Erd\H{o}s and G. N. S\'ark\"ozy,
\emph{On cycles in the coprime graph of integers},
Electronic Journal of Combinatorics \textbf{4}(2) (1997), R8.
\href{https://doi.org/10.37236/1323}{doi:10.37236/1323}.

\bibitem{S99}
G. N. S\'ark\"ozy,
\emph{Complete tripartite subgraphs in the coprime graph of integers},
Discrete Mathematics \textbf{202} (1999), 227--238.
\href{https://doi.org/10.1016/S0012-365X(98)00359-8}{doi:10.1016/S0012-365X(98)00359-8}.

\bibitem{Bloom}
T. F. Bloom, \emph{Erd\H{o}s Problem 883}, original statement and historical references,
\url{https://www.erdosproblems.com/883}.

\bibitem{DP26}
D. Della Pietra, \emph{Erd\H{o}s 883 asymptotic proof manuscript},
2026 repository snapshot
\texttt{1f26276d7850fc9afe68e9d7edecf24aa7cc597d}.
\url{https://github.com/donalddellapietra/erdos-883-proof/blob/1f26276d7850fc9afe68e9d7edecf24aa7cc597d/paper/main.tex}.

\bibitem{Pan26}
Y. Pan, \emph{A Lean verified finite completion of the odd cycle assertion in Erd\H{o}s Problem 883},
manuscript dated 5 October 2026, Theorem 1 and supporting source.
\url{https://github.com/zoahdev/erdos883-first-question/blob/44a32ef7faead567b2961cb0ec8a1859abf461f5/paper/source/manuscript.tex}.
Cited as related proof work, not as an organizer-verified priority determination.

\bibitem{Archive}
C. Jentges, AI-assisted project development,
\emph{Odd cycles in the coprime graph of dense sets: Erd\H{o}s Problem 883(i) for every $n$},
mathematical source and certificate archive, 10 October 2026.
Repository \texttt{wizardscurtain/contest}, at the immutable commit and directory
specified in Section~\ref{sec:evidence}. The public companion release must include
\sourcepath{evidence/hall_tail_20261010} and its manifest dependencies.

\bibitem{TailReplay}
C. Jentges, AI-assisted project record,
\emph{Fresh analytic-tail replay and mathematical supplement},
11 October 2026. Files \sourcepath{ANALYTIC_TAIL_REPLAY.md} and
\sourcepath{analytic_interval_replay.json}, at the source commit and supplement
directory specified in Section~\ref{sec:evidence}. The reported execution covers
the analytic tail, not a fresh all-$n$ verification.
\end{thebibliography}

\end{document}
